\chapter{What a Quantum Field Is}

This chapter follows Leonard Susskind's third lecture in \emph{Particle Physics 1: Basic Concepts}, curated by LazyingArt LLC. We begin exactly where the lecture begins: last time we were beginning to get to the idea of a quantum field, and now we want to answer two linked questions. What is a quantum field, and how is it related to particles? The lecture does not answer that by postulate. It rebuilds the answer from periodic space, harmonic oscillators, occupation numbers, and Fourier modes, and only then allows the field itself to emerge.

\section{Periodic Space, Allowed Wave Numbers, and Momentum}

We begin, not with operators, but with periodic space. The circle is a trick, but it is a very useful trick. If we put a system in a finite interval with hard walls, waves reflect, and momentum is no longer the right organizing principle. If instead we identify the ends of the interval, a wave can move around the circle without reflecting. That is not merely an analog of momentum conservation; for this finite box it is the way to preserve momentum conservation while avoiding the complications of infinite volume.

A plane wave on a circle of circumference \(L\) has the form \(e^{ikx}\). Periodicity requires that after traveling once around the circle the wave come back to its original value:
\begin{equation}
e^{ikL}=1.
\end{equation}
Since \(e^{i2\pi n}=1\) for every integer \(n\), the allowed wave numbers satisfy
\begin{equation}
kL=2\pi n,
\qquad
k=\frac{2\pi n}{L},
\qquad
n\in\mathbb{Z}.
\end{equation}
The quantity \(k\) is the wave number. It is important to keep track of what is and is not an integer here. The integer is \(n\). The wave number \(k\) itself is an integer multiple of \(2\pi/L\), not an integer.

The same condition can be stated in terms of wavelength. If \(\lambda\) is the distance for one full cycle, then one full cycle means \(k\lambda=2\pi\), so
\begin{equation}
k\lambda=2\pi,
\qquad
k=\frac{2\pi}{\lambda}.
\end{equation}
Thus large wavelength means small wave number, and small wavelength means large wave number. The lecture then reconnects this to ordinary quantum mechanics: for a quantum-mechanical wave, momentum and wave number are related by
\begin{equation}
p=\hbar k.
\end{equation}
In units with \(\hbar=1\), one often speaks loosely of \(k\) as the momentum. That is the lecture's practice too, but only after the distinction has been made clearly.

We should also keep one more point in view, because the late questions return to it. Making space periodic does not discretize position. What becomes discrete is the set of allowed momenta. The positions remain points on the circle. Later, when \(L\) is made larger and larger, the spacing \(2\pi/L\) between neighboring allowed values of \(k\) becomes smaller and smaller, and the discrete sum over modes tends toward an integral.

\section{Occupation Numbers and Fock-Space State Labels}

Once the allowed momenta are in place, the lecture asks the next natural question: how many quanta occupy each allowed momentum mode? Here Susskind slows down, because the notation can become misleading if we are careless. One integer, \(n\), labels the allowed values \(k=2\pi n/L\). A different integer tells us how many quanta sit in a given mode. The whole transition to field language depends on not confusing these two roles.

The second integer is the occupation number, written \(n(k)\). It means the number of particles, or quanta, with wave number \(k\). A many-particle state is therefore specified by giving the full list of occupation numbers:
\begin{equation}
\big|\{n(k)\}\big\rangle
=
\big|\ldots,n(-2),n(-1),n(0),n(1),n(2),\ldots\big\rangle.
\end{equation}
The lecture first writes this schematically as a string of occupations for the first, second, third, and so on allowed momenta, and then reminds us that the string must extend in both directions, because both positive and negative momenta occur.

\begin{figure}[t]
\centering
\includegraphics[width=0.72\textwidth]{lecture_03_figure_02.png}
\caption{Occupation-number notation \(n(k)\) on the board. The board shows the bookkeeping device; the text writes the state cleanly as a full list of occupations.}
\end{figure}

The vacuum is simply the state in which every slot is empty:
\begin{equation}
|0\rangle
=
|\ldots,0,0,0,0,0,\ldots\rangle.
\end{equation}
At this stage the lecture is very concrete. Nothing has yet been said about interactions, or even about the field itself. We only have a set of allowed momentum modes and a way to specify how many quanta occupy each of them. But this is already a genuine many-particle state space, and it is the right one for quantum fields.

\section{Mode Operators \(a^\dagger(k)\) and \(a(k)\)}

Now that the states have been labeled, the next step is to introduce operators that act on those labels. On the board Susskind writes \(a^{+}(k)\) and \(a^{-}(k)\). We will use the standard notation \(a^\dagger(k)\) and \(a(k)\), but the meaning is the same: for each momentum mode \(k\), there is one operator that raises the occupation number and one that lowers it.

Their action on an occupation-number state is
\begin{align}
a^\dagger(k)\big|\ldots,n(k),\ldots\big\rangle
&=
\sqrt{n(k)+1}\,\big|\ldots,n(k)+1,\ldots\big\rangle,
\\
a(k)\big|\ldots,n(k),\ldots\big\rangle
&=
\sqrt{n(k)}\,\big|\ldots,n(k)-1,\ldots\big\rangle.
\end{align}
The lecture briefly stumbles over the wording and then corrects itself in exactly the right way: \(a^\dagger(k)\) goes to the \(k\)-th slot and increases that occupation number by one unit, while \(a(k)\) decreases it, bringing in the appropriate square root.

\begin{figure}[t]
\centering
\includegraphics[width=0.58\textwidth]{lecture_03_figure_03.png}
\caption{Mode-by-mode creation and annihilation operators on the board. The lecture writes \(a^{+}(k)\) and \(a^{-}(k)\); the notes use \(a^\dagger(k)\) and \(a(k)\).}
\end{figure}

A one-particle momentum eigenstate is obtained by acting on the vacuum with a single creation operator:
\begin{equation}
|k\rangle = a^\dagger(k)|0\rangle.
\end{equation}
With this notation, the lecture immediately shows how the operator language can describe scattering. If a particle initially in mode \(k_7\) is scattered into mode \(k_9\), then
\begin{equation}
a^\dagger(k_9)a(k_7)|k_7\rangle = |k_9\rangle.
\end{equation}
Why do the square roots not show up here? Because \(a(k_7)\) acts on a slot whose occupation number is \(1\), giving \(\sqrt{1}=1\), and \(a^\dagger(k_9)\) creates into an empty slot, again giving \(\sqrt{0+1}=1\).

The same language can also describe processes in which particle number changes. That is the point at which single-particle quantum mechanics really runs out of room. Schematically,
\begin{equation}
a^\dagger(k_9)a^\dagger(k_{16})a(k_7)|k_7\rangle
=
|k_9,k_{16}\rangle.
\end{equation}
One particle goes in, two come out. The lecture is careful to say that this is not yet a law of nature; it is a mathematical description of a possible kind of process. If the dynamics allows particle number to change, this is the language in which such changes are naturally expressed.

\section{From Classical Fourier Coefficients to the Quantum Field}

Only now does the lecture ask the decisive question: what is the connection between the field and the \(a\)'s? The answer is given by moving briefly back to the classical picture. A classical field on a periodic interval can be expanded in Fourier modes,
\begin{align}
\psi(x) &= \sum_k \alpha(k)e^{ikx},
\\
\psi^\ast(x) &= \sum_k \alpha^\ast(k)e^{-ikx}.
\end{align}
The coefficients \(\alpha(k)\) are just the Fourier amplitudes of the classical wave. Since the real and imaginary parts of those amplitudes can be measured, the lecture treats them as classical observables.

The quantum step is then stated quite directly: the \(a\)'s are the quantum-mechanical versions of the Fourier coefficients. The lecture does not write an ordinary equality sign here, because these are not the same kinds of objects. Instead, it treats quantization as replacement:
\begin{equation}
\alpha(k)\rightsquigarrow a(k),
\qquad
\alpha^\ast(k)\rightsquigarrow a^\dagger(k).
\end{equation}
Susskind openly hesitates for a moment about which convention is standard and then flips to the standard choice. That hesitation is pedagogically useful. It reminds us that the sign conventions are conventions. In the rest of the chapter we fix
\begin{align}
\Psi(x) &= \sum_k a(k)e^{ikx},
\\
\Psi^\dagger(x) &= \sum_k a^\dagger(k)e^{-ikx}.
\end{align}
This is the point at which the field finally appears: a quantum field is an operator-valued Fourier expansion.

There is also an important warning at this stage. Creation and annihilation operators are mathematical operators. Their existence does not mean that the laws of nature must create and destroy particles. Whether particle number changes is a dynamical question. But if the laws of nature do permit such processes, then these operators are the right tools for describing them. The lecture refuses to let the terminology do more work than the mathematics has earned.

\subsection{Question \& Answer}

What does \(\Psi^\dagger(x)\) actually create?

This is the cleanest conceptual obstacle in the lecture, and it is resolved by a direct calculation. First, the annihilation field kills the vacuum:
\begin{equation}
\Psi(x)|0\rangle = 0.
\end{equation}
There is no particle present to annihilate. The creation field is more interesting. Acting on the vacuum, it gives
\begin{align}
\Psi^\dagger(x)|0\rangle
&=
\sum_k e^{-ikx}a^\dagger(k)|0\rangle
\\
&=
\sum_k e^{-ikx}|k\rangle,
\end{align}
where
\begin{equation}
|k\rangle=a^\dagger(k)|0\rangle.
\end{equation}

\begin{figure}[t]
\centering
\includegraphics[width=0.70\textwidth]{lecture_03_figure_05.png}
\caption{\(\Psi^\dagger(x)\) acting on the vacuum. The board first expands the field in momentum modes and then rewrites the result directly in the one-particle momentum basis.}
\end{figure}

At first this looks like a strange object: a sum of momentum eigenstates with phases \(e^{-ikx}\). But that is exactly the construction of a position state in one-particle quantum mechanics. If we superpose momentum states with those phases, we get a state localized at position \(x\). So the lecture answers its own question in the natural way:
\begin{equation}
\Psi^\dagger(x)\ \text{creates a particle at position }x.
\end{equation}
Equivalently, \(\Psi(x)\) annihilates a particle at position \(x\).

There is a secondary observation here which the lecture also makes. Since each coefficient \(e^{-ikx}\) has unit magnitude, all allowed momenta appear with equal probability weight in this state. So the position state is a coherent superposition spread uniformly over momentum space. If we enlarge the circle and let \(L\to\infty\), the sum over \(k\) will become an integral, but the interpretation does not change.

\section{Local Reactions, Bose Enhancement, and Stimulated Emission}

Once \(\Psi(x)\) and \(\Psi^\dagger(x)\) have acquired a position-space meaning, the lecture immediately turns to locality. This is where the tone becomes particularly revealing. Susskind first proposes a ridiculous process on purpose: a particle is absorbed at a sharply localized point \(x\), and then two particles are emitted somewhere far away, one on Alpha Centauri and one on Betelgeuse. That, he says, is not the sort of process we want. It is manifestly non-local.

So what is the right local process? If one particle is absorbed at a point \(x\) and two are emitted, then the emission should occur at the same point. That is the motivation for the operator product
\begin{equation}
\Psi^\dagger(x)\Psi^\dagger(x)\Psi(x).
\end{equation}
The first field annihilates the incoming particle at \(x\), and the two creation fields emit the outgoing particles from that same point.

\begin{figure}[t]
\centering
\includegraphics[width=0.92\textwidth]{lecture_03_figure_06.png}
\caption{The board layout during the locality and probability discussion. On the left are the field expansions; on the right is the momentum-space amplitude for the local \(1\to 2\) process.}
\end{figure}

\begin{equation}
\Psi(x)=\sum_q a(q)e^{iqx},
\qquad
\Psi^\dagger(x)=\sum_\ell a^\dagger(\ell)e^{-i\ell x}.
\end{equation}
These are precisely the formulas carried on the left board in the lecture, rewritten here in a consistent convention.

\begin{figure}[t]
\centering
\begin{tikzpicture}[scale=1.0]
\draw[->] (-2.2,0) -- (0,0);
\fill (0,0) circle (2.2pt);
\draw[->] (0,0) -- (2.0,1.1);
\draw[->] (0,0) -- (2.0,-1.1);
\node at (-1.2,0.28) {$k_{\mathrm{initial}}$};
\node[right] at (0.05,0.05) {$x$};
\node at (1.45,1.35) {$\ell$};
\node at (1.45,-1.35) {$m$};
\end{tikzpicture}
\caption{Editorial redraw of the local process emphasized in the lecture: one incoming particle is absorbed at the point \(x\), and two outgoing particles are emitted from that same point.}
\end{figure}

Let the initial state be a one-particle momentum eigenstate \(|k_{\mathrm{initial}}\rangle\). The annihilation field acts first:
\begin{align}
\Psi(x)|k_{\mathrm{initial}}\rangle
&=
\sum_q e^{iqx}a(q)a^\dagger(k_{\mathrm{initial}})|0\rangle
\\
&=
e^{ik_{\mathrm{initial}}x}|0\rangle.
\end{align}
Only the term matching the occupied mode survives. The lecture stresses this point: an annihilation operator has to find something to annihilate.

Now the two creation fields act on the vacuum:
\begin{align}
\Psi^\dagger(x)\Psi^\dagger(x)\Psi(x)|k_{\mathrm{initial}}\rangle
&=
e^{ik_{\mathrm{initial}}x}
\sum_{\ell,m}
e^{-i\ell x}e^{-imx}
a^\dagger(\ell)a^\dagger(m)|0\rangle
\\
&=
\sum_{\ell,m}
e^{i(k_{\mathrm{initial}}-\ell-m)x}
|\ell,m\rangle.
\end{align}
This is the basic local amplitude. The final state is a superposition of all two-particle momentum states. In this simplified model we have deliberately allowed the target to absorb recoil, so there is no exact momentum conservation among the particles alone. The phase records the difference between the incoming momentum and the two outgoing momenta, which is exactly the momentum carried away by the target.

At first the lecture draws the conclusion that all final momentum pairs appear with equal probability, since the coefficient is only a phase. Then it pauses and corrects itself. The special case \(\ell=m\) is different. If the first creation operator has already created one particle in mode \(m\), then the second creation into that same mode produces an extra factor:
\begin{equation}
a^\dagger(m)\big|\ldots,n(m)=1,\ldots\big\rangle
=
\sqrt{2}\,\big|\ldots,n(m)=2,\ldots\big\rangle.
\end{equation}
So there is a Bose enhancement when the two particles come out in the same state.

This is not the end of the story, however. A student points out that when \(\ell\neq m\), the raw double sum contains both ordered terms
\begin{equation}
a^\dagger(\ell)a^\dagger(m)|0\rangle
\qquad\text{and}\qquad
a^\dagger(m)a^\dagger(\ell)|0\rangle,
\end{equation}
and these are equal because creation operators commute. So the lecture backs away from an overly strong comparison between the same-momentum and different-momentum channels. The safe conclusion, and the one worth preserving, is the algebraic one: whenever we create into a mode that is already occupied, an extra square-root factor appears. That is the bosonic effect we should trust.

The lecture then gives the cleaner example. Suppose an atom at the point \(x\) emits a photon. If the initial state is the vacuum, then
\begin{equation}
\Psi^\dagger(x)|0\rangle
=
\sum_k e^{-ikx}|k\rangle,
\end{equation}
so in this simplified model every momentum mode is emitted with equal magnitude. That is the spontaneous process.

Now let there already be one photon in mode \(\ell\). Then
\begin{equation}
\Psi^\dagger(x)|\ell\rangle
=
\sum_{k\neq \ell} e^{-ikx}|k,\ell\rangle
+
e^{-i\ell x}\sqrt{2}\,|2_\ell\rangle,
\end{equation}
where \(|2_\ell\rangle\) denotes the state with two quanta in the same mode \(\ell\). The extra \(\sqrt{2}\) is exactly the enhancement seen in the local \(1\to 2\) calculation.

More generally, if the mode \(\ell\) already has occupation number \(n_\ell\), then
\begin{equation}
a^\dagger(\ell)\big|\ldots,n_\ell,\ldots\big\rangle
=
\sqrt{n_\ell+1}\,
\big|\ldots,n_\ell+1,\ldots\big\rangle.
\end{equation}
Therefore the probability for emission into that already occupied mode is enhanced by a factor
\begin{equation}
n_\ell+1.
\end{equation}
If \(n_\ell=0\), this is spontaneous emission. If \(n_\ell\) is large, the enhancement can be very large. That is stimulated emission.

The lecture is careful not to dress this up in the wrong language. Bosons do not literally attract one another into the same state. The effect is statistical. It comes from the identity of the particles and the square-root factors in the bosonic algebra.

A final technical remark from the lecture belongs here. In more complicated operator products, the order is not always obvious at first sight. But because creation operators commute among themselves and annihilation operators commute among themselves, one can rearrange expressions into normal order, with annihilation operators to the right and creation operators to the left. In that ordering, one first lets the operators eat the incoming particles and then create the outgoing ones.

\section{Bosonic Algebra, Number Density, and Late Clarifications}

Only after these examples does the lecture permit itself to summarize the algebra. This order matters. The commutation relations are not introduced as abstract axioms floating above the physics; they appear after we have already seen what they do.

For bosons,
\begin{align}
[a^\dagger(k),a^\dagger(\ell)] &= 0,
\\
[a(k),a(\ell)] &= 0,
\\
[a(k),a^\dagger(\ell)] &= \delta_{k\ell}.
\end{align}
Here \(\delta_{k\ell}\) is the Kronecker delta, equal to \(1\) when \(k=\ell\) and \(0\) otherwise. This is the harmonic-oscillator algebra, repeated for each momentum mode.

The lecture then asks what object is represented by the integral
\begin{equation}
\int_0^L dx\,\Psi^\dagger(x)\Psi(x).
\end{equation}
Rather than naming it immediately, it works it out:
\begin{align}
\int_0^L dx\,\Psi^\dagger(x)\Psi(x)
&=
\sum_{k,\ell}
a^\dagger(k)a(\ell)
\int_0^L dx\,e^{-ikx}e^{i\ell x}
\\
&=
\sum_{k,\ell}
a^\dagger(k)a(\ell)
\int_0^L dx\,e^{i(\ell-k)x}.
\end{align}
Now the periodicity of the circle does the work once again:
\begin{equation}
\int_0^L dx\,e^{i(\ell-k)x}=L\,\delta_{k\ell}.
\end{equation}
So the double sum collapses:
\begin{equation}
\int_0^L dx\,\Psi^\dagger(x)\Psi(x)
=
L\sum_k a^\dagger(k)a(k).
\end{equation}
But \(a^\dagger(k)a(k)\) is just the occupation number in the mode \(k\). Hence, with the lecture's normalization convention,
\begin{equation}
\frac{1}{L}\int_0^L dx\,\Psi^\dagger(x)\Psi(x)
=
\sum_k a^\dagger(k)a(k)
=
N,
\end{equation}
the total number of particles. The lecture explicitly remarks that other conventions absorb the factor of \(L\) into the field. Here we leave the factor visible, because that is how the lecture keeps track of it.

Accordingly, \(\Psi^\dagger(x)\Psi(x)\) is the particle-density operator. One might be tempted to call it the number of particles at the point \(x\), but density is the better language.

\subsection{Question \& Answer}

Why is the vacuum not the zero vector?

This question is important because the lecture had earlier said that annihilating the vacuum gives zero. That sentence is correct, but only if one distinguishes very carefully between the vacuum state and the null vector.

The vacuum \(|0\rangle\) is a genuine state. It is the ground state of the system, the state with no particles in it, and it is normalized:
\begin{equation}
\langle 0|0\rangle=1.
\end{equation}
The zero vector, by contrast, is not a physical state at all. It is the null element of the vector space.

So when we write
\begin{equation}
a(k)|0\rangle=0,
\end{equation}
the \(0\) on the right-hand side is the zero vector. The statement is that annihilating the vacuum produces no state. It is not the statement that the vacuum itself is the zero vector. The lecture insists on this distinction, and then pushes it one step further: the field is not a state either. In quantum field theory the field is an operator, and the states are the things on which it acts.

Several late clarifications now fall into place. In more than one spatial dimension, the periodic box becomes a torus. In two dimensions it is a rectangle with opposite edges identified; in three dimensions it is a cube with opposite faces identified. The point of the construction is still the same as in one dimension: one makes space periodic in order to avoid infinities while keeping momentum as the natural organizing variable. What becomes discrete is momentum, not position.

The lecture also revisits the analogy with the harmonic oscillator. The real and imaginary parts of the field are the analogs of the oscillator's conjugate variables. Up to conventional factors,
\begin{equation}
\Re\Psi \propto \Psi+\Psi^\dagger,
\qquad
\Im\Psi \propto \frac{\Psi-\Psi^\dagger}{i}.
\end{equation}
These are conjugate quantities, and they do not commute. So one cannot measure all components of the field simultaneously with arbitrary precision. But when the occupation numbers are very large, the amplitudes are large compared with the fixed right-hand side of the commutator. In that regime the field behaves approximately classically. This is why a large number of photons can be described, to very good approximation, by an ordinary classical electromagnetic wave. It is also why an external classical field can be viewed, in the appropriate limit, as a state containing very many quanta.

Finally the lecture previews the fermionic alternative. The key algebraic change is that commutators are replaced by anticommutators:
\begin{equation}
\{A,B\}=AB+BA.
\end{equation}
In the bosonic theory, many particles can pile into the same state and produce Bose enhancement. In the fermionic theory the Pauli exclusion principle forbids that. Schematically, the bosonic relation
\begin{equation}
[a(k),a^\dagger(\ell)]=\delta_{k\ell}
\end{equation}
is replaced by an anticommutator relation of the same general form, and the result is the opposite statistical behavior: no stimulated emission into an already occupied fermionic state, and no classical field limit obtained by piling indefinitely many quanta into one mode.

\section{Summary}

The lecture has built the simplest bosonic quantum field theory in the order in which one would want to understand it. We begin with a field on a circle, quantize the allowed wave numbers, and connect \(k\), \(\lambda\), and momentum. We then distinguish sharply between the integer that labels an allowed momentum mode and the occupation number \(n(k)\) that counts how many quanta sit in that mode.

From there the harmonic-oscillator operators \(a^\dagger(k)\) and \(a(k)\) act mode by mode on the Fock-space state labels. The major conceptual step is then to reinterpret those operators as quantized Fourier coefficients. That is what turns the classical wave \(\psi(x)\) into the operator-valued field \(\Psi(x)\). Acting with \(\Psi^\dagger(x)\) on the vacuum produces the familiar momentum superposition that one-particle quantum mechanics recognizes as a position state, so the field becomes an operator that creates and annihilates particles at points in space.

Only after that positional meaning is established does the lecture insist on locality, write the local reaction \(\Psi^\dagger(x)\Psi^\dagger(x)\Psi(x)\), and let the bosonic square-root factors show their physical content through Bose enhancement and stimulated emission. At the end the algebra is summarized, the density operator \(\Psi^\dagger(x)\Psi(x)\) is identified, and the late clarifications return us to the larger picture: the vacuum is not the zero vector, periodic space is a finite-volume device rather than a lattice ontology, large occupation numbers recover classical field behavior, and fermions will require a different algebra altogether.